
\[
  \chi(F\otimes L)=\chi(F)
\]
\item
\emph{\(L\) est num\'eriquement \'equivalent \`a \(0\).}
\end{enumerate}
\emph{Si \(X/k\) est projectif et muni d'un module inversible ample \(H={O}_X(1)\), les conditions pr\'ec\'edentes sont aussi \'equivalentes aux suivantes}~:
\begin{enumerate}[label=(\roman*),start=4]
\item \emph{Pour tout entier \(m\) (positif ou n\'egatif), \(L^{\otimes m}(1)\) est ample.}
\item \emph{(Si \(X\) est irr\'eductible). Pour tout couple d'entiers \(m,n\), on a}
\[
  \chi(L^{\otimes m}(n))=\chi({O}_X(n))
\]
\item \emph{(Si \(X\) est irr\'eductible de dimension \(r\ge2\)). On a}
\[
  \langle c_1(L)\cdot c_1(H)^{r-1}\rangle=0,\qquad
  \langle c_1(L)^2\cdot c_1(H)^{r-2}\rangle=0.
\]
\end{enumerate}

D\'emontrons que (i) implique (ii). Appliquant le lemme de d\'evissage (EGA III 3.1) et 4.5 (i), on se ram\`ene au cas o\`u \(X\) est int\`egre et \(F={O}_X\). Alors on doit prouver que le polyn\^ome \(\chi(L^{\otimes m})\) est constant. Pour ceci, on peut remplacer \(L\) par un multiple, donc supposer \(L\) alg\'ebriquement \'equivalent \`a \(0\). Dans ce cas, si \(L\) est d\'eform\'e en \({O}\) par la famille alg\'ebrique d\'efinie par un faisceau inversible \(M\) sur un \(X\times_kT\) o\`u \(T\) est connexe, de type fini, le polyn\^ome \(\chi(M_t^{\otimes m})\) ne d\'epend pas de \(t\in T\) (EGA III 7.9.5)\,; d'o\`u l'implication.

Soit \(Y\) une courbe int\`egre ferm\'ee dans \(X\). Faisant \(F={O}_Y\) dans (ii), on trouve \(\chi({O}_Y)=\chi(L_Y)\). Donc la formule (``th\'eor\`eme de Riemann'') \(\chi(L_Y)=\langle c_1(L)\cdot Y\rangle+\chi({O}_Y)\) donne \(\langle c_1(L)\cdot Y\rangle=0\)\,; d'o\`u (ii) implique (iii).

Soit \(f:X'\to X\) une pr\'esentation de Chow~: \(X'\) est projectif et \(f\) est surjectif. Si l'on admet que (iii) implique (i) pour \(X'\), alors 
